\documentclass{article}
\input{../../../lib.sty}

\title{\texttt{fn new\_conditional\_rdp\_curve}}
\author{Michael Shoemate}

\begin{document}
\maketitle

Proves soundness of \texttt{fn new\_conditional\_rdp\_curve} in \texttt{mod.rs}.

\section{Hoare Triple}
\subsection*{Precondition}
\texttt{base\_curve} upper-bounds the base Rényi divergence at every order $> 1$, and \texttt{x}
is such that \texttt{log\_inv\_gamma\_hi\_from\_x(x)} succeeds whenever it is called.

\subsection*{Pseudocode}
\lstinputlisting[language=Python,firstline=2,escapechar=|]{./pseudocode/new_conditional_rdp_curve.py}

\subsection*{Postcondition}
\begin{theorem}
\label{postcondition}
For every $\alpha > 1$, the returned function either returns \texttt{Err(e)} due to a propagated
helper error, or returns an upper bound on the Corollary 16 expression
\[
\epsilon(\alpha)
\;+\;
\frac{\alpha - 2}{\alpha - 1}\epsilon(\alpha - 1)
\;+\;
\frac{2 \log(1/\gamma)}{\alpha - 1},
\]
where $\epsilon(\cdot)$ is the base Rényi divergence curve and $\gamma = 1 - e^{-x}$.
\end{theorem}

\section{Proof}
Fix $\alpha > 1$.

\begin{lemma}
\label{guard}
If $\alpha \le 2$, the returned value is $+\infty$, which is a valid upper bound.
\end{lemma}

\begin{proof}
The expression is finite or infinite, and $+\infty$ upper-bounds either case.
\end{proof}

\begin{lemma}
\label{main}
If $\alpha > 2$ and the function returns \texttt{Ok(out)}, then \texttt{out} upper-bounds the
displayed Corollary 16 expression.
\end{lemma}

\begin{proof}
Let $\epsilon(\cdot)$ denote the exact base Rényi divergence curve.
By precondition,
\[
\texttt{eps\_alpha} \ge \epsilon(\alpha).
\]
Also \texttt{alpha\_minus\_one\_hi} $\ge \alpha - 1 > 1$.
Rényi divergence is nondecreasing in its order, so
\[
\epsilon(\alpha - 1) \le \epsilon(\texttt{alpha\_minus\_one\_hi})
\le \texttt{eps\_alpha\_minus\_one\_hi}.
\]
Because $\alpha > 2$, both $\alpha - 2$ and $\alpha - 1$ are positive. The arithmetic helpers
therefore give
\[
\texttt{coeff\_hi} \ge \frac{\alpha - 2}{\alpha - 1}.
\]
Since \texttt{eps\_alpha\_minus\_one\_hi} is nonnegative, upward-rounded multiplication yields
\[
\texttt{term1\_hi} \ge \frac{\alpha - 2}{\alpha - 1}\epsilon(\alpha - 1).
\]
By the proof of \texttt{log\_inv\_gamma\_hi\_from\_x},
\[
\texttt{log\_inv\_gamma\_hi} \ge \log(1/\gamma).
\]
Again \texttt{alpha\_minus\_one\_lo} $\le \alpha - 1$ with a positive denominator, so
\[
\texttt{term2\_hi} \ge \frac{2 \log(1/\gamma)}{\alpha - 1}.
\]
Finally, \texttt{out} is formed by upward-rounded addition of these three upper bounds, so it
upper-bounds the full displayed expression.
\end{proof}

\begin{proof}[Proof of Theorem~\ref{postcondition}]
By Lemma~\ref{guard} for $\alpha \le 2$, and Lemma~\ref{main} otherwise.
\end{proof}

\end{document}
